\documentclass[letterpaper,11pt]{article}
\usepackage[top=0.8in, bottom=0.8in, left=0.9in, right=0.9in]{geometry}
\usepackage[hidelinks]{hyperref}
\usepackage{lmodern}
\usepackage{parskip}
\usepackage{microtype}
\setlength{\parskip}{6pt}
\setlength{\parindent}{0pt}
\begin{document}
\begin{flushleft}
\textbf{\large Ajay Venkatesh} \\
Brooklyn, NY \\
+1 516 585 9013 \\
\href{mailto:av3855@nyu.edu}{av3855@nyu.edu}
\end{flushleft}
\vspace{10pt}
May 21, 2026
\vspace{6pt}

Hiring Manager \\
Mercury Insurance

\vspace{10pt}

Dear Hiring Manager,

My name is Ajay Venkatesh. I'm finishing my M.S. in Computer Engineering at NYU, with production software experience at PwC in Bengaluru, a summer SDE internship at Amazon, and ongoing on-campus work at NYU's Novel AI Technologies. I'm writing about the MLOps Catastrophe Modeling and Geo-Analytics role.

The combination, productionizing ML models with high availability and reliability through automated pipelines, is exactly the work I've been doing. At Amazon I deployed a \textbf{Retrieval-Augmented Generation (RAG)} pipeline with an \textbf{agentic workflow} that autonomously reasons over operational logs to generate context-aware summaries, and implemented a \textbf{Model Context Protocol (MCP)} server that materially reduced query time and compute costs. I provisioned resilient \textbf{Infrastructure-as-Code} on \textbf{AWS CDK} across auto-scaling ECS clusters with multi-stage CI/CD and rollback controls, the same posture MLOps deployment demands.

On modeling and training: my Big Data project built a \textbf{PySpark / HDFS} pipeline over millions of flight records, training Random Forest and Gradient Boosted Tree models with rigorous EDA. I've also worked on \textbf{RoBERTa fine-tuning with custom LoRA} and \textbf{ResNet} architectures under strict parameter and memory budgets, and built RAG systems on \textbf{AWS Bedrock} with FAISS indexing.

On data pipelines and engineering: at NYU I engineered \textbf{ETL pipelines} transforming \textbf{NoSQL} into \textbf{PostgreSQL} with concurrency control, processing hundreds of thousands of records and powering \textbf{AWS QuickSight} dashboards. At PwC I parallelized batch ingestion on operational feeds that materially cut execution time. Reliable data plumbing feeding production models is the unglamorous part of MLOps that compounds.

Stack-wise: \textbf{Python}, \textbf{SQL}, \textbf{PyTorch}, \textbf{TensorFlow}, \textbf{MLflow}, plus \textbf{AWS} (Lambda, SageMaker, ECS Fargate, CDK, S3, IAM), \textbf{Docker}, \textbf{Kubernetes}, \textbf{Git}, and \textbf{GitHub} CI/CD workflows. Comfortable with cloud infrastructure for high-availability ML systems.

What pulls me toward Mercury is the operational seriousness. Catastrophe modeling for insurance runs on accuracy that translates into real-world risk pricing and claim outcomes, where careful MLOps engineering matters more than novelty.

I'd welcome the chance to discuss further. Thanks for considering my application.

\vspace{12pt}
Sincerely, \\
Ajay Venkatesh

\end{document}